\section{Conclusion}
\label{sec:conclusion}

This survey has told one continuous story, in the order it actually happened. It runs from Ashley's
HYPO in 1991, through case-based reasoning's reframing as lazy learning in 2003, through the
statistical and neural prediction of legal outcomes between 2014 and 2018, through domain-adapted
legal language models and their benchmarks between 2020 and 2023, through the emergence of legal
retrieval-augmented generation in 2024, through a year of maturing systems and sharper evaluation in
2025, and on to the contemporary systems of 2026, Legalyze among them. If we read the field this
way, instead of organizing it by topic, a specific pattern becomes visible, one a topic-first
account can obscure: each stage responded to a concrete limitation of the one before it. Case-based
reasoning showed that legal relevance depends on more than textual similarity, but it struggled to
scale. Quantitative outcome prediction scaled, but could not explain its own reasoning. Legal
language models processed text fluently, but could not guarantee they were right.
Retrieval-augmented generation connected generation to external evidence, but retrieval quality,
legal correctness, and citation support turned out to be three separate problems, not one.

The comparative case study of Legalyze and Ren et al.\ (2026) shows this same pattern at the scale
of a single year. Two systems appeared in the same year, addressed the same general problem, and
chose to solve different halves of it. One used structure-aware evidence grounding; the other used
synthetic training data. Neither choice is complete on its own. That is exactly the lesson this
survey's reliability synthesis drew from \citet{butler2026}'s error taxonomy: a hallucination, a
retrieval error, and a reasoning error call for three different fixes, and a system that only
addresses one of them will still fail on the other two. The jurisdiction and multilingual synthesis
extends this lesson geographically. In contexts such as Pakistan, effective systems have to connect
user queries, which users may express in English, Urdu, or Roman Urdu, to the appropriate legal
sources while preserving statutory hierarchy, judicial context, and provenance. A system that
researchers validate only on English-language, single-jurisdiction data has not yet shown that it
generalizes.

Future research should prioritize structure-preserving retrieval, provenance-aware legal knowledge
graphs, multilingual query processing, claim-level citation verification, and agentic workflows with
explicit quality gates. Researchers should test these through reproducible benchmarks that cover the
complete evidence-to-answer pipeline and report retrieval effectiveness, legal correctness, citation
support, and cost, without collapsing them into a single score, in the spirit of the groundedness,
retrieval-accuracy, and correctness decomposition this survey has described. The goal of legal AI
should not be to replace legal judgment with fluent automated responses. It should be to make legal
information more accessible, systematic, and auditable, while keeping the distinction between
retrieved evidence, model-generated interpretation, and conclusions that need professional review.
Thirty-five years after Ashley first asked whether a computer could compare legal cases, the field's
next stage will depend less on generation alone and more on the ability to retrieve, contextualize,
verify, and reproduce evidence-supported legal reasoning.
